\section{The model, its closures and the formal results}\label{app:proofs}

\noindent
This appendix collects the formal apparatus behind the closure matrix: the
notation and the functional forms that solve the producer and household
problems of Section~\ref{subsec:optimization}
(Section~\ref{app:model-kernel}); the five labor-market closures and the
three financing closures (Sections~\ref{app:labor-closures}
and~\ref{app:financing-closures}); and the formal results --- the closure
systems, the demand-only reduction and the rigidity result
(Proposition~\ref{prop:nest};
Sections~\ref{app:closure-systems}--\ref{app:sectoral}). The
intermediate-bill and external-account accounting is in
Appendix~\ref{app:abill}.

\subsection{Production and consumption}\label{app:model-kernel}

\paragraph{Notation.} $L_i$ is employment in sector $i$ and $L=\sum_i L_i$ total employment; a bar marks baseline employment, so $\bar L_i$ is baseline sectoral and $\Lbar$ baseline total employment. $w$ is the common wage of the mobile closures, $w_i$ the
sector-specific wage where labor markets are sectoral. Real wages are computed by deflating with the consumption price index $\mathcal{P}(p)$, the CES price of the household's bundle; a $0$ marks all
other baseline values. $L_i^{\mathrm{cm}}(p,y_i,w_i)$ is the labor demand that
solves the firm's cost-minimization problem
of Section~\ref{subsec:optimization} at prices $p$ and output $y_i$. The
allocation parameter $\eta$ takes only its endpoints: $\eta=1$ lets labor
move to its cost-minimizing use, $\eta=0$ freezes the baseline allocation of labor.
The labor-supply elasticity $\eta_s$ governs how strongly sectoral employment
responds to the real wage; in principle, it may differ across sectors ($\eta_{s,i}$), but the paper uses a
common value $\eta_s$ throughout. Finally, $\theta$, $\varepsilon$ and $\sigma$ are the
substitution elasticities of Section~\ref{subsec:optimization} (among
intermediates, between labor and intermediates, and in consumption); their
limit $0^{+}$ is fixed proportions. Table~\ref{tab:parameters} collects these elasticities and the two parameters related to the extensions in Section~\ref{sec:beyond-demand-only} for convenience.

% =====================================================================
% Table --- the elasticities of the comparison (Sec. 4.2)
% ---------------------------------------------------------------------
% Requested by reviewer 1 (R1.3: "including elasticity values in a table").
% Values are read from the executed design experiments/designs/
% matrix_5x3.toml and from the manifests of the reported cells
% (runs/matrix_5x3-*): sigma, theta, epsilon, eta are the design values;
% eta_s is the bracket printed in Table 4; the two extension parameters are
% the ones the reported tables of Section 6 use. None of these is estimated
% from the input--output table.
% =====================================================================
\begin{table}[t]
\centering
\small
\begin{tabular}{@{}l p{5.9cm} l l@{}}
\toprule
Parameter & Meaning in the model core & Value & Set in \\
\midrule
$\sigma$ & substitution across consumption goods & $0.9$ & all but the Leontief column \\
$\theta$ & substitution across intermediates & $0.5$ & all but the Leontief column \\
$\varepsilon$ & labor--intermediate nesting & $0.5$ & all but the Leontief column \\
$\eta_s$ & sectoral labor-supply elasticity with respect to the real wage & $\{0.5,1,2\}$ & elastic-supply column \\
\addlinespace


$\eta$ & intersectoral labor-allocation selector & $\{0,1\}$ & labor closure \\
$\delta$ & capacity extension: utilization elasticity of productivity & $\{-0.5,0,+0.5\}$ & Section~\ref{subsec:pricing-capacity} \\
$\kappa$ & markup extension: margin--utilization slope & $\{0.25,0.5\}$ & Section~\ref{subsec:pricing-markup} \\

\bottomrule
\end{tabular}
\par\smallskip
\parbox{\textwidth}{\footnotesize
\textit{Note:} $\eta=0$ fixes employment in each sector at its baseline level,
with sector-specific wages; $\eta=1$ selects cost-minimizing labor demand,
with a common wage in the standard CGE and wage-pinned closures and sectoral
wages under elastic supply. The three substitution elasticities take the
corner values $(\theta,\varepsilon,\sigma) = (10^{-4},10^{-4},10^{-4})$ in the
Leontief column, the fixed-coefficient corner solution of
Appendix~\ref{app:proofs}; the welfare measures of
Table~\ref{tab:welfare-comparison} are evaluated at the reference
$\sigma = 0.9$ in every column.}
\caption{Parameter values for the baseline comparison and the pricing extensions. The first block lists the substitution elasticities of the model core --- $\sigma$ across consumption goods, $\theta$ across intermediates, $\varepsilon$ in the labor--intermediate nesting --- and the sectoral labor-supply elasticity $\eta_s$. The second block lists the labor-allocation selector $\eta$ and the two extension parameters $\delta$ (capacity) and $\kappa$ (markup) at the values executed. The last column states where each value is set; none of these parameters is estimated from the input--output table.}
\label{tab:parameters}
\end{table}


Each sector's technology is a nested CES with constant returns: labor and an intermediate composite substitute with elasticity $\varepsilon$, and intermediate inputs substitute with elasticity $\theta$ inside the composite. Cost minimization, Equation~\eqref{eq:cost-min}, gives the intermediate price index and the unit cost
\begin{align}
P_i(p) &= \Big(\textstyle\sum_{u=1}^N \Omega_{iu}\,p_u^{1-\theta}\Big)^{1/(1-\theta)},
\label{eq:model-ip}\\
c_i(p,w_i;A_i^{H}) &= \Big((A_i^{H})^{\varepsilon-1}\big(\alpha_i\,w_i^{1-\varepsilon}
  + (1-\alpha_i)\,P_i^{1-\varepsilon}\big)\Big)^{1/(1-\varepsilon)},
\label{eq:model-cost}
\end{align}
with the Cobb--Douglas limits at $\theta=1$ and $\varepsilon=1$. Because
$\sum_u\Omega_{iu}=1$, the unit cost equals one at the calibration point
$(p,w,A^{H})=(\one,\one,\one)$. Zero profit requires
\begin{equation}
p_i = c_i(p,w_i;A_i^{H}) \qquad \forall\, i.
\label{eq:zp}
\end{equation}
Shephard's lemma, evaluated at zero profit ($p_i = c_i$), gives the
cost-minimizing labor demand
\begin{equation}
L_i^{\mathrm{cm}}(p,y_i,w_i)
 = \alpha_i (A_i^{H})^{\varepsilon-1} y_i
   \left(\frac{p_i}{w_i}\right)^{\varepsilon},
\label{eq:closure-labor-demand}
\end{equation}
where $y_i$ is gross output, $A_i^{H}$ Hicks-neutral productivity and $\alpha_i$ the calibrated
value-added share. Shephard's lemma gives demand for the whole intermediate
composite~\eqref{eq:model-ip}, with share $1-\alpha_i$; domestic demand is its
fraction $a_i/(1-\alpha_i)$, which replaces $1-\alpha_i$ by $a_i$
(see \eqref{eq:int}).

The household has CES preferences with elasticity $\sigma$ and calibrated
weights $\beta_i$. The consumption price index, which is the numeraire in the flexible-wage closures (the wage-pinned closures normalize $w=1$ instead), is
\begin{equation}
  \mathcal{P}(p) = \left(\sum_{i=1}^N \beta_i p_i^{1-\sigma}\right)^{1/(1-\sigma)},
  \qquad \sum_{i=1}^N \beta_i = 1,
  \label{eq:labor-cpi}
\end{equation}
with the Cobb--Douglas limit at $\sigma = 1$. Given household consumption
expenditure $E^h$, utility maximization gives
\begin{equation}
c_i^h = E^h\frac{\beta_i p_i^{-\sigma}}{\sum_j\beta_j p_j^{1-\sigma}}.
\label{eq:model-household}
\end{equation}
Goods-market clearing for each of the $N$ sectors reads
\begin{equation}
y_i = \mathit{int}_i + (1-m_i)\big(c_i^h + g_i^G + \mathit{inv}_i + g_i\big) + X_i,
\label{eq:model-clearing}
\end{equation}
where $m_i$ is the sectoral import margin, $g^G$ baseline government
purchases, $\mathit{inv}$ non-program investment, $X$ exports and $g$ the
program bundle (zero under F1). Every closure below enforces all $N$
clearings~\eqref{eq:model-clearing}.

\subsection{The five labor closures}\label{app:labor-closures}

All closures use the same labor demand~\eqref{eq:closure-labor-demand}; they
differ in the wage and labor-supply restrictions imposed on it.

\paragraph{The no-reallocation closure: immobile labor.}
The no-reallocation column fixes employment in every sector and solves for a sector-specific wage to clear each labor market:
\begin{equation}
L_i = L_i^{\mathrm{cm}}(p,y_i,w_i) = \bar{L}_i,
\qquad i=1,\ldots,N, \qquad \eta = 0.
\label{eq:closure-bf}
\end{equation}
Thus $\sum_i L_i = \Lbar$, but wages may differ across sectors. Setting $\eta = 1$ instead gives the standard CGE closure below.

\paragraph{The standard CGE closure: mobile labor at full employment.}
Labor reallocates across sectors at one flexible wage, while total
employment remains fixed:
\begin{equation}
L_i = L_i^{\mathrm{cm}}(p,y_i,w),
\qquad \sum_{i=1}^N L_i = \Lbar,
\qquad \eta = 1.
\label{eq:closure-alpha}
\end{equation}

\paragraph{The elastic labor-supply closure: sectoral supply curves.}
In the elastic-supply column, each sector has its own wage and labor-supply curve, and all sectors share one supply elasticity $\eta_s$:
\begin{equation}
L_i = L_i^{\mathrm{cm}}(p,y_i,w_i)
 = \bar{L}_i \left[\frac{w_i/\mathcal{P}(p)}{w_i^{0}/\mathcal{P}_0}\right]^{\eta_s},
\qquad \eta_s \geq 0, \qquad i=1,\ldots,N.
\label{eq:closure-beta}
\end{equation}
Labor supply responds to the real wage. The deflator is the consumption price index $\mathcal{P}(p)$, which is the same for all sectors, not the sector's own price $p_i$. The baseline real wage $w_i^{0}/\mathcal{P}_0$ is calibrated to 1. The supply curves are reduced forms and have no income effect on the labor--leisure choice. As Section~\ref{sec:two-questions} explains, $\eta_s$ is a labor-supply elasticity, not the allocation parameter $\eta$. Even so, setting $\eta_s = 0$ reproduces the no-reallocation column exactly.

The results matrix (Table~\ref{tab:results-matrix}) and the welfare comparison (Table~\ref{tab:welfare-comparison}) report the elastic-supply column as a range over $\eta_s \in [0.5, 2]$. The figures and the pricing tables use $\eta_s = 0.5$. Table~\ref{tab:sectoral} adds further values of $\eta_s$ under F2.

\paragraph{The employment-adjustment closure: wage pinned, employment uncapped.}
No labor-market equation is imposed, so employment is demand-determined and the money wage serves as the numeraire:
\begin{equation}
w = 1 \quad(\text{numeraire}),
\qquad F = 0,
\qquad L_i = L_i^{\mathrm{cm}}(p,y_i,w),
\qquad \eta = 1 .
\label{eq:closure-gamma}
\end{equation}
Unlike the standard CGE closure, this closure does not impose $\sum_i L_i = \Lbar$: employment absorbs whatever labor equilibrium production requires.

\paragraph{The Leontief closure: the fixed-coefficient corner.}
The Leontief closure keeps the employment rule of the preceding closure while taking the substitution elasticities to their Leontief limits:
\begin{equation}
w = 1, \qquad F = 0, \qquad \eta = 1,
\qquad (\theta,\varepsilon,\sigma) \longrightarrow (0^+,0^+,0^+).
\label{eq:closure-delta}
\end{equation}
It is the fixed-coefficient special case of the employment-adjustment closure and has no labor-market mechanism of its own. The two columns have the same solution as long as prices stay at their baseline values, which is the case at the calibration technology under a demand-only impulse (Lemma~\ref{lem:affine}). Once prices move, the solutions differ (Section~\ref{sec:beyond-demand-only}).

\subsection{The three financing closures}\label{app:financing-closures}

Each labor closure is crossed with three financing specifications calibrated from the same reference program. Let $g_i = G_0\psi_i$ denote the reference program bundle, where $G_0$ is its baseline-price value and $\psi_i \geq 0$ with $\sum_i\psi_i = 1$ describes its sectoral incidence; a $k$-fold program is $kg$. F2 and F3 add this bundle to final demand; F1 instead uses its scale and incidence to redirect household spending.

Baseline government purchases $g_i^G$ remain tax-financed in all three cases. Define the wage bill and their current-price tax cost as 
\begin{equation}
  W = \sum_{i=1}^N w_i L_i, \qquad T_G(p) = \sum_{i=1}^N p_i g_i^G. \label{eq:financing-baseline-budget} 
\end{equation}
With a common wage, $W = w\sum_i L_i$. Let $E$ denote after-tax household resources, including the net external transfer $F$ from the rest of the world,\footnote{$F$ is the variable that clears the external account; the labels F1--F3 name financing closures, not a variable.} and let $s$ be the calibrated saving rate. Household consumption expenditure is $E^h = (1-s)E$. Gross household purchases $c_i^h$ exhaust this budget:
\begin{equation} 
\sum_{i=1}^N p_i c_i^h = E^h, \qquad c_i^{\mathrm{d}} = (1-m_i)c_i^h. \label{eq:financing-household-budget}
\end{equation}
Only the domestic content enters domestic goods-market clearing~\eqref{eq:model-clearing}; likewise, the program contributes $(1-m_i)g_i$ under F2 and F3.

\paragraph{F1: demand repurposing.}
F1 changes the household's CES weights without adding a public-demand
bundle. For positive preference shifters $d_i$, the weights and gross
household demand become
\begin{align}
\widetilde{\beta}_i
 &= \frac{\beta_i d_i}{\sum_j\beta_j d_j},
\qquad \sum_i\widetilde{\beta}_i = 1,
\label{eq:closure-f1-weights}\\
c_i^h
 &= E^h\frac{\widetilde{\beta}_i p_i^{-\sigma}}
 {\sum_j\widetilde{\beta}_j p_j^{1-\sigma}},
\qquad E = W-T_G(p)+F.
\label{eq:closure-f1}
\end{align}
The preference shifter is $d_i = 1+kG_0\psi_i^+/c_{i0}^h$ in sectors with positive baseline household purchases, and $d_i=1$ otherwise; $\psi^+$ is the program incidence restricted to those sectors and renormalized to sum to one. The CES normalizer preserves \eqref{eq:financing-household-budget} at every price vector. For a given household budget, higher spending shares in program sectors are therefore offset by lower shares elsewhere, although equilibrium income can still change. F1 models a shift in household demand, caused for example by changed preferences or by regulation. It adds no autonomous spending, so it is not an investment injection, and its effects should not be read as a multiplier.

\paragraph{F2: tax-financed public investment.}
F2 adds the real program bundle $g$ and levies an additional household
tax that covers the program at equilibrium prices:
\begin{equation}
T(p) = \sum_{i=1}^N p_i g_i,
\qquad E = W-T_G(p)-T(p)+F,
\qquad B_{\mathrm{gov}}=0.
\label{eq:closure-f2}
\end{equation}
Total domestic tax revenue $T_G(p)+T(p)$ therefore balances the purchases
$g^G+g$, priced at current rather than baseline prices. The tax is
implemented proportionally as $\tau=[T_G(p)+T(p)]/W$ and the tax base moves with the wage bill $W$.

\paragraph{F3: externally financed public investment.}
F3 adds the same real bundle but imposes no additional domestic program
tax. Its current-price cost is booked as external program financing:
\begin{equation}
B_{\mathrm{gov}} = \sum_{i=1}^N p_i g_i,
\qquad E = W-T_G(p)+F.
\label{eq:closure-f3}
\end{equation}
The booking $B_{\mathrm{gov}}$ is distinct from the net external transfer
$F$ entering household resources after tax, although their economic interpretation is similar, namely some form of external debt. $F$ is solved endogenously in the no-reallocation, standard CGE and elastic-supply closures, but is fixed at zero in the two wage-pinned closures, the employment-adjustment closure and its Leontief corner. $B_{\mathrm{gov}}$, by contrast, is set endogenously to match the size of the
spending program.
\begin{equation}
S + T_{\mathrm{int}} + M - (I+X) = F+B_{\mathrm{gov}},
\qquad S=sE,
\label{eq:financing-external-account}
\end{equation}
where $I$ is non-program investment, $X$ exports, $M$ total imports (including intermediate and program imports), and $T_{\mathrm{int}}$ the separately booked product-tax leak on intermediate use; all are current-price values. F3 closes the open-economy accounts without specifying how external financing occurs: because the model is static, it treats borrowing and transfers identically and omits interest payments, repayment, and exchange-rate adjustment.
\subsection{The closure systems}\label{app:closure-systems}

We keep the notation above and record the data-calibrated objects the
systems below use: domestic bill coefficients $a_i$, the value of domestic
intermediates used per unit of gross output of sector $i$, the input-share matrix $\Omega$ ($\Omega_{iu}$ the share of input $u$ in
sector $i$'s bundle, rows summing to one), sector import margins $m_i$, the
saving rate $s$, baseline government demand $g^G$, non-program investment
$inv$, exports $X$, and the labor bar $\Lbar = \sum_i \bar L_i = 1$ in
baseline income units (GDP $=1$).

\paragraph{Price block.} The price block is \eqref{eq:model-ip}--\eqref{eq:zp}; at
$\varepsilon = 1$ the unit cost takes the Cobb--Douglas form
$c_i = w^{\alpha_i} P_i^{\,1-\alpha_i}/A_i^{H}$.

\paragraph{Quantity block.} With full cost-minimizing allocation
($\eta = 1$) the marginal-product condition inverts to the labor demand
\eqref{eq:closure-labor-demand}. At $\eta = 0$ the allocation is the
frozen baseline, $L_i = \bar{L}_i$; under the sectoral-wage closure the wage
$w$ in \eqref{eq:closure-labor-demand} is then a per-sector $w_i$, solved from
$L^{cm}_i(p, y_i, w_i) = \bar{L}_i$ jointly with \eqref{eq:zp}.
Intermediate demand is charged with the \emph{domestic}
bill:
\begin{equation}
int_i = p_i^{-\theta} \sum_j \Omega_{ji}\,
        p_j^{\varepsilon} (A_j^{H})^{\varepsilon-1} P_j^{\theta-\varepsilon}\, a_j\, y_j.
\label{eq:int}
\end{equation}
Domestic household demand, $c_i^{\mathrm{d}}=(1-m_i)c_i^h$, is CES with CPI
deflation and the weights $\widetilde{\beta}$ of \eqref{eq:closure-f1-weights}
($\widetilde{\beta} = \beta$ under F2 and F3):
\begin{equation}
c_i^{\mathrm{d}} = (1-s)(1-m_i)\,\widetilde{\beta}_i\; E\; p_i^{-\sigma} \Big/
      \textstyle\sum_j \widetilde{\beta}_j\, p_j^{1-\sigma},
\qquad
E = (1-\tau)\, W + F,
\label{eq:cons}
\end{equation}
with the wage-income tax rate $\tau = p\!\cdot\! g^G / W$ under
F1/F3 and $\tau = p\!\cdot\!(g^G + g)/W$ under F2, and $F$ the
net external transfer entering expenditure after tax. Saving
$sE$ leaks; the import content is charged margin by margin. Goods-market
clearing is \eqref{eq:model-clearing}, with $g_i$ present only under the
additive financings F2/F3 ($g_i = 0$ under F1).

\paragraph{The four regimes.} The standard CGE closure (mobile labor) solves for
$(p, y, w, F) \in \mathbb{R}^{2N+2}$ with the consumption price index
\eqref{eq:labor-cpi} as numeraire and the labor equation
\eqref{eq:closure-alpha} closing the system. The sectoral elastic-supply closure
\eqref{eq:closure-beta} replaces the single money wage by $N$ sectoral wages and
is solved in $(p, y, w(1{:}N), F) \in \mathbb{R}^{3N+1}$, with its $N$ supply
curves closing the system at the anchors $w_i^{0}/\mathcal{P}_0 = 1$. The wage-pinned systems solve for $(p, y) \in \mathbb{R}^{2N}$ with the
wage pinned $w \equiv 1$ and $F \equiv 0$; employment is computed from the
solution. The employment-adjustment closure uses the baseline elasticities,
while the Leontief closure sets them to the Leontief
corner $(\theta,\varepsilon,\sigma) = (10^{-4},10^{-4},10^{-4})$. The no-reallocation closure is the sectoral-wage system at
$\eta = 0$: the
allocation stays frozen at the baseline, $L_i = \bar{L}_i$, and $N$ sectoral
wages $w_i$ are solved from each sector's own first-order condition
$L^{cm}_i(p, y_i, w_i) = \bar{L}_i$ jointly with zero profit, in the same unknowns as the
elastic-supply closure and with the external transfer $F$ free. All columns
in the matrix use $\eta = 1$ except the no-reallocation closure.

\subsection{The demand-only reduction}\label{app:reduction}

\begin{lemma}[Price block demand-invariance]\label{lem:price}
The zero-profit system \eqref{eq:zp} is homogeneous of degree one in
$(p,w)$ jointly and contains no demand variable. Consequently, for any
technology $A^{H}$, every solution satisfies $p = w\,\pi(A^{H})$, where
$\pi(A^{H}) \in \mathbb{R}^N_{++}$ solves $c_i(\pi,1;A^{H}) = \pi_i$ for all $i$;
relative prices and the ratio $p/w$ are functions of $A^{H}$ alone. Under the
CPI numeraire, $w = 1/\mathcal{P}(\pi(A^{H}))$ and $p = \pi(A^{H})/\mathcal{P}(\pi(A^{H}))$; under the
wage pin, $w = 1$ and $p = \pi(A^{H})$. At $A^{H} \equiv \one$ the unit cost equals one
at the calibration point, so $\pi = \one$, hence in both gauges
\begin{equation}
p = \one, \qquad w = 1, \qquad \mathcal{P} = 1,
\end{equation}
independently of the demand vector, of the financing closure, and of the
shock size. The lemma concerns the \emph{single-wage} price block; with
sectoral wages it does not apply (Appendix~\ref{app:sectoral}). The statement is scoped to the positive price solution under consideration:
$\pi(A^{H})$ denotes that solution at $w=1$,
and the equivalences below compare the systems on that branch.
\end{lemma}

\begin{proof}
Homogeneity: $c_i(tp, tw; A^{H}) = t\, c_i(p,w;A^{H})$ for $t>0$, by inspection of
\eqref{eq:model-cost}--\eqref{eq:model-ip}. Write $q = p/w$. Then \eqref{eq:zp} is
equivalent to $q_i = c_i(q,1;A^{H})$ for all $i$, i.e.\ $q = \pi(A^{H})$, and
the corresponding solutions of \eqref{eq:zp} form the ray
$\{(w\,\pi(A^{H}), w):w>0\}$. The
CPI \eqref{eq:labor-cpi} is homogeneous of degree one in $p$, so
$\mathcal{P}(w\,\pi) = w\,\mathcal{P}(\pi)$; the numeraire $\mathcal{P} = 1$ or the pin $w = 1$
selects the scale. At $A^{H} \equiv \one$, $\pi = \one$ solves
\eqref{eq:zp} at $w = 1$ because $c_i(\one,1;\one) = 1$, and $\mathcal{P}(\one) = 1$ because $\sum_i \beta_i = 1$.
\end{proof}

\begin{corollary}[Real-wage invariance]\label{cor:realwage}
For the standard CGE, employment-adjustment and Leontief closures,
every solution of the shared block satisfies
\begin{equation}
\frac{w}{\mathcal{P}(p)} = \frac{1}{\mathcal{P}(\pi(A^{H}))},
\end{equation}
a function of technology alone. At $A^{H} \equiv \one$ this equals $1 =
w_0/\mathcal{P}_0$: the equilibrium real wage sits at its anchor whatever the
demand.
\end{corollary}

\begin{lemma}[Elasticity-free affine quantity block]\label{lem:affine}
Fix $\eta = 1$ and $A^{H} \equiv \one$, and evaluate the quantity block at the
demand-only prices $(p,w) = (\one,1)$. Define
\[
\varphi = (1-s)\,\mathrm{diag}(1-m)\,\widetilde{\beta},
\]
with $\widetilde{\beta}$ the household weights of \eqref{eq:cons}.
For every substitution-elasticity triple $(\theta,\varepsilon,\sigma)$,
goods-market clearing \eqref{eq:model-clearing} is exactly affine,
\begin{equation}
y = G\, y + b(F),
\qquad
G = \Omega^{\!\top}\mathrm{diag}(a)
  + \varphi\,\alpha^{\!\top},
\label{eq:affine}
\end{equation}
\begin{equation}
b(F) = (1-m)\odot\big( g^G + inv + g \big) + X
     + \varphi
       \big( F - \textstyle\sum g^G - \mathbf{1}_{F2}\textstyle\sum g \big),
\label{eq:bvec}
\end{equation}
where $g$ is present only under F2/F3 and $\mathbf{1}_{F2}$ selects the
tax-financed row.
\end{lemma}

\begin{proof}
At these prices all price and productivity factors in \eqref{eq:int} and
\eqref{eq:closure-labor-demand} equal one, giving
$int = \Omega^{\!\top}\mathrm{diag}(a)y$ and $L_i = \alpha_i y_i$,
while \eqref{eq:cons} becomes $c^{\mathrm{d}} = \varphi E$.
Substituting $E = \alpha^{\!\top}y - \sum g^G - \mathbf{1}_{F2}\sum g + F$
into \eqref{eq:model-clearing} gives \eqref{eq:affine}--\eqref{eq:bvec}.
\end{proof}

\begin{remark}[Nonsubstitution]\label{rem:samuelson}
Lemmas~\ref{lem:price} and~\ref{lem:affine} are the production side of the
nonsubstitution theorem \citep{Samuelson.1951}: with one primary factor,
constant returns and no joint production, relative prices are independent of
demand, so the fixed-price multiplier \citep{Robinson.2006} is exact at the
calibration technology for every substitution-elasticity triple.
\end{remark}

\paragraph{What the transfer does.} In the mobile regime the labor equation
selects the transfer. In \eqref{eq:bvec}, $b(0)$ is autonomous
demand at zero transfer for the chosen financing rule, including the program
where present and the tax-induced reduction in household demand. Then
$b(F) = b(0) + \varphi F$, where $\varphi$ is the domestic consumption demand
per unit of transfer. Writing $y(F) = (I-G)^{-1}(b(0) + \varphi F)$, the pin
$\alpha^{\!\top}y = \Lbar$ gives
\begin{equation}
F^{\ast} = \frac{\Lbar - \alpha^{\!\top}(I-G)^{-1} b(0)}
               {\alpha^{\!\top}(I-G)^{-1} \varphi}.
\label{eq:finstrument}
\end{equation}
This instrument is why the F2 and F3 rows within the mobile-labor column are real-neutral: F2
charges the program through the wage-income tax \eqref{eq:closure-f2} ($E = \alpha^{\!\top}y -
\sum g^G - \sum g + F$), F3 through the external account \eqref{eq:closure-f3} ($E = \alpha^{\!\top}y
- \sum g^G + F$), and the two $b$-vectors differ by exactly
$\varphi\,\sum g$. The instrument \eqref{eq:finstrument} therefore yields
$F^{\ast}_{F3} = F^{\ast}_{F2} - \sum g$, identical real allocations and an
identical net external position $F + B_{\mathrm{gov}}$ (measured:
$F_{F3} - F_{F2} = -G_0$ in every mobile column). Where no labor equation remains,
nothing determines $F$, which is why the wage-pinned columns fix it at zero.

\subsection{The elastic labor supply model and its two ends}\label{app:sectoral}

The elastic-supply column is one member of a family indexed by $\eta_s$. At $\eta_s=0$ the family coincides with the no-reallocation column (Proposition~\ref{prop:nest}); as $\eta_s\to\infty$ it converges to the employment-adjustment allocation, but at a different external transfer.

\paragraph{The model.} The elastic-supply form \eqref{eq:closure-beta} gives
each sector its own wage and supply curve; the matrix's elastic-supply
column is this system for $\eta_s \in [0.5, 2]$, and
Table~\ref{tab:sectoral} reports results for a wider range of $\eta_s$ under F2, including both ends. Because
each sector's own wage enters its unit cost $c_i(p,w_i;A_i^{H})$ directly, while the
intermediate price index \eqref{eq:model-ip} contains goods prices only, the $N$
zero-profit conditions now depend on the wage structure and through it on the
employment levels: Lemma~\ref{lem:price} no longer applies, and demand reaches
prices through the wages, $p = \pi(A^{H}, w)$.

\begin{proposition}[The rigid corner is the $\eta = 0$ endpoint]\label{prop:nest}
Set $\eta_s = 0$. Then the elastic labor supply model and the $\eta = 0$
endpoint of the matrix have the same solution set.
\end{proposition}

\begin{proof}
Write both systems in the unknowns $(p, y, w, F)$ as four row blocks:
\begin{align*}
\text{(P)} \quad & p_i = c_i(p, w_i; A^{H}), & & i = 1,\dots,N,
   & & \text{\eqref{eq:zp}} \\
\text{(Y)} \quad & y_i = int_i + c_i^{\mathrm{d}} + (1-m_i)(g^G_i + inv_i + g_i) + X_i,
   & & i = 1,\dots,N, & & \text{\eqref{eq:model-clearing}} \\
\text{(W)} \quad & \log L^{cm}_i(p,y_i,w_i) - \log \bar{L}_i
   - \eta_s \log \frac{w_i}{\mathcal{P}(p)} = 0, & & i = 1,\dots,N, & & \\
\text{(E)} \quad & \mathcal{P}(p) = 1. & & & &
\end{align*}
The quantity block (Y) draws labor income from the wage bill
$\sum_i w_i L_i$ inside \eqref{eq:cons} and from the tax rate $\tau$; this is
the only place, besides (W), where the two systems differ. The $\eta = 0$
endpoint is the same four blocks with two changes: (W) drops the supply term
($\eta_s = 0$), and (Y) evaluates the wage bill at the \emph{frozen}
allocation, $L_i := \bar{L}_i$, instead of at the cost-minimizing demand
$L^{cm}_i$.

\emph{(i)} Let $(p,y,w,F)$ solve the sectoral system at $\eta_s \equiv 0$. Its (W)
gives $L^{cm}_i(p,y_i,w_i) = \bar{L}_i$ for every $i$; substituting this into (Y) makes
that block equal the endpoint's (Y), while (P) and (E) are identical in the two
systems by inspection, so $(p,y,w,F)$ solves the endpoint system.
\emph{(ii)} Conversely, let $(p,y,w,F)$ solve the endpoint system. Its (W) gives
$L^{cm}_i(p,y_i,w_i) = \bar{L}_i$, which satisfies the sectoral (W) at
$\eta_s \equiv 0$, and substituting it into (Y) yields the sectoral (Y).
Both systems are square --- $3N+1$ unknowns, $3N+1$ equations --- and their
residual maps agree on the common solution set, so the solution sets coincide.
\end{proof}

As the common elasticity $\eta_s$ is raised beyond the values used in
Table~\ref{tab:results-matrix}, employment and consumption converge and the
transfer turns positive. The limit does not coincide with the
employment-adjustment closure, because the elastic-supply model determines the
external transfer $F$ endogenously while the employment-adjustment closure sets
$F = 0$. Solved at the transfer selected
by the elastic-supply model, $F_{1000}$ (the transfer at $\eta_s=1000$), the wage-pinned system
reproduces that model's allocation (employment $1.002324$ against $1.002327$,
consumption $-1.289$ against $-1.288$) and keeps prices exactly at the
baseline, whereas the elastic-supply prices deviate from it by
$8.6 \times 10^{-5}$, the residual from approximating the limit by a finite
$\eta_s$. At $F = 0$ the two allocations differ by the demand effect of the
transfer; Table~\ref{tab:sectoral} reports the limit at both transfers.

% Table B1 -- the sectoral family in three parts (rigid end / uniform ladder /
% elastic end).  Every cell is a model solution under F2 financing, demand-only
% shocks, and baseline elasticities theta = eps = 0.5, sigma = 0.9, eta = 1,
% varying only the sectoral labor-supply elasticity eta_s; each row is
% reproducible by re-solving the corresponding sectoral-family specification.
\begin{table}[htbp]
\centering
\small
\begin{tabular}{lrrrr}
\toprule
& $\max|p-1|$ & Employment & Consumption & $F$ \\
\midrule
\multicolumn{5}{@{}l}{\emph{Uniform elasticity}} \\
$\eta_s = 0.25$ & $0.152897$ & $1.000864$ & $-1.689$ & $-0.004424$ \\
$\eta_s = 0.5$  & $0.105666$ & $1.001255$ & $-1.571$ & $-0.002828$ \\
$\eta_s = 1$    & $0.065422$ & $1.001627$ & $-1.467$ & $-0.001444$ \\
$\eta_s = 2$    & $0.037171$ & $1.001914$ & $-1.391$ & $-0.000457$ \\
$\eta_s = 10$   & $0.008351$ & $1.002231$ & $-1.312$ & $+0.000562$ \\
$\eta_s = 100$  & $0.000859$ & $1.002318$ & $-1.291$ & $+0.000829$ \\
\midrule
\multicolumn{5}{@{}l}{\emph{The rigid end}} \\
Sectoral, $\eta_s = 0$    & $0.279009$ & $1.000000$ & $-1.980$ & $-0.008528$ \\
No reallocation ($\eta = 0$)  & $0.279009$ & $1.000000$ & $-1.980$ & $-0.008528$ \\
\midrule
\multicolumn{5}{@{}l}{\emph{The elastic end}} \\
Sectoral, $\eta_s = 1000$      & $0.000086$ & $1.002327$ & $-1.288$ & $+0.000856$ \\
Employment adjustment, $F = 0$      & $0.000000$ & $1.001260$ & $-1.533$ & $0.000000$ \\
Employment adjustment, $F = F_{1000}$ & $0.000000$ & $1.002324$ & $-1.289$ & $+0.000856$ \\
\bottomrule
\end{tabular}
\caption{Sectoral labor supply at different elasticities, all under financing closure F2. Price deviations are fractions; employment is an index with baseline of one; consumption is a percentage change; $F$ is expressed as a share of baseline GDP. \emph{Top:} a common
elasticity across sectors, ranging from $\eta_s = 0.25$ to $100$. \emph{Middle:} the rigid end, where the family at $\eta_s = 0$ and the no-reallocation row coincide. \emph{Bottom:} a large-elasticity approximation at $\eta_s = 1000$, compared with the employment-adjustment system at $F=0$ and at the sectoral economy's transfer $F_{1000}$. At the same external transfer the two closures give nearly equal employment and consumption responses.}
\label{tab:sectoral}
\end{table}


Two further points follow. First, solving the sectoral system at $\eta_s=0$
from the $\eta = 0$ solution reproduces all four blocks exactly, so no separate
rigid sectoral cell is reported; this case coincides with the no-reallocation
column of Table~\ref{tab:results-matrix}. Second, the proof of
Proposition~\ref{prop:nest} does not use $A^{H}\equiv\one$, so the nesting also
holds under supply shocks. Under such shocks the elastic end separates and
$\eta_s$ becomes identifiable; under a demand-only program it remains a
sensitivity parameter.
